\documentclass[conference]{IEEEtran}
\IEEEoverridecommandlockouts

\usepackage{cite}
\usepackage{amsmath,amssymb,amsfonts}
\usepackage{algorithmic}
\usepackage{graphicx}
\usepackage{textcomp}
\usepackage{xcolor}
\usepackage{bm}
\usepackage{url}

\def\BibTeX{{\rm B\kern-.05em{\sc i\kern-.025em b}\kern-.08em
    T\kern-.1667em\lower.7ex\hbox{E}\kern-.125emX}}

\begin{document}

\title{D-EPNet: An EP-DIFFormer Detector for Massive MIMO}

\author{
\IEEEauthorblockN{Author Name}
\IEEEauthorblockA{
Affiliation \\
City, Country \\
email@domain.com
}
}

\maketitle

\begin{abstract}
This paper presents D-EPNet, a model-driven detector for uplink massive multiple-input multiple-output (MIMO) systems. The proposed method combines expectation propagation (EP) with a DIFFormer-based graph refinement network. In each iteration, EP produces symbol-wise Gaussian cavity statistics, and a neural refinement module further exploits channel-induced symbol coupling through linear-time diffusion attention, graph aggregation, and gated recurrent state updates. Compared with a conventional EP detector, the proposed architecture preserves the probabilistic inference structure while introducing a more expressive interaction model for symbol refinement. The resulting detector outputs symbol logits over the pulse-amplitude modulation alphabet and is trained by cross-entropy loss.
\end{abstract}

\begin{IEEEkeywords}
massive MIMO, expectation propagation, DIFFormer, graph neural network, signal detection
\end{IEEEkeywords}

\section{Introduction}

Massive MIMO detection can be formulated as an inference problem over a high-dimensional discrete symbol vector. Classical model-based algorithms exploit the linear system structure, whereas neural detectors can learn complicated interference patterns from data. The implementation provided in this work adopts a hybrid strategy: an EP module produces tractable Gaussian cavity distributions, and a DIFFormer-based graph network refines these beliefs by using both graph-aware propagation and global interaction modeling.

It is worth noting that the Python wrapper keeps the legacy class name \texttt{ChebNet} for compatibility, but the actual neural refinement module is no longer a Chebyshev graph network. Instead, the current implementation uses the class \texttt{sp\_GNN\_Difformer}, whose core message-passing block is \texttt{DIFFormer\_v2}. The following description is written according to the current code rather than the older Chebyshev-based formulation.

\section{System Model}

Consider an uplink massive MIMO system with $N_t$ single-antenna users and a base station equipped with $N_r$ receive antennas. Let $\mathbf{x}_c \in \mathbb{C}^{N_t}$ denote the transmitted symbol vector and let $\mathbf{y}_c \in \mathbb{C}^{N_r}$ denote the received signal vector. The complex baseband model is
\begin{equation}
\mathbf{y}_c = \mathbf{H}_c \mathbf{x}_c + \mathbf{n}_c,
\label{eq:complex_model}
\end{equation}
where $\mathbf{H}_c \in \mathbb{C}^{N_r \times N_t}$ is the channel matrix and $\mathbf{n}_c \sim \mathcal{CN}(\mathbf{0}, \sigma_c^2 \mathbf{I})$ is additive white Gaussian noise.

Following the implementation, we convert \eqref{eq:complex_model} into an equivalent real-valued model:
\begin{equation}
\mathbf{y}=
\begin{bmatrix}
\Re(\mathbf{y}_c)\\
\Im(\mathbf{y}_c)
\end{bmatrix},\quad
\mathbf{x}=
\begin{bmatrix}
\Re(\mathbf{x}_c)\\
\Im(\mathbf{x}_c)
\end{bmatrix},\quad
\mathbf{n}=
\begin{bmatrix}
\Re(\mathbf{n}_c)\\
\Im(\mathbf{n}_c)
\end{bmatrix},
\end{equation}
\begin{equation}
\mathbf{H}=
\begin{bmatrix}
\Re(\mathbf{H}_c) & -\Im(\mathbf{H}_c)\\
\Im(\mathbf{H}_c) & \Re(\mathbf{H}_c)
\end{bmatrix}.
\end{equation}
The resulting system is
\begin{equation}
\mathbf{y} = \mathbf{H}\mathbf{x} + \mathbf{n},
\label{eq:real_model}
\end{equation}
where $\mathbf{H}\in\mathbb{R}^{N\times K}$ with $N=2N_r$, $K=2N_t$, and $\mathbf{n}\sim \mathcal{N}(\mathbf{0}, \sigma_n^2 \mathbf{I})$ with $\sigma_n^2=\sigma_c^2/2$.

For $4^k$-QAM modulation, each real-valued symbol belongs to the $2^k$-PAM alphabet
\begin{equation}
\mathcal{X} = \{\pm 1, \pm 3, \ldots, \pm (2^k-1)\}.
\label{eq:alphabet}
\end{equation}
Hence, detection can be viewed as estimating $\mathbf{x}\in\mathcal{X}^K$ from $\mathbf{y}$ and $\mathbf{H}$, or equivalently as a $|\mathcal{X}|$-class classification problem for each entry of $\mathbf{x}$.

The Bayesian posterior can be written as
\begin{equation}
p(\mathbf{x}|\mathbf{y},\mathbf{H})
\propto
\exp\!\left(
-\frac{1}{2\sigma_n^2}\|\mathbf{y}-\mathbf{H}\mathbf{x}\|_2^2
\right)
\prod_{k=1}^{K} p(x_k),
\label{eq:posterior}
\end{equation}
where $p(x_k)$ is uniform over $\mathcal{X}$.

\subsection{Graph Preprocessing}

The training script first generates a graph-oriented preprocessing output through
\begin{equation}
(\mathbf{L}_{\mathrm{cep}}, \hat{\mathbf{y}}, \alpha) = \Gamma(\mathbf{H}, \mathbf{y}),
\label{eq:preprocess}
\end{equation}
which corresponds to the function \texttt{sp\_GNN\_data(H, y)} in the code. Here, $\mathbf{L}_{\mathrm{cep}}\in\mathbb{R}^{K\times K}$ is a dense symbol-interaction matrix, $\hat{\mathbf{y}}\in\mathbb{R}^K$ is a node-wise observation statistic, and $\alpha$ is a channel-dependent scaling factor.

For the $k$-th symbol, the initial node feature used by the neural module is
\begin{equation}
\mathbf{s}_k^{(0)}=
\left[
\hat{y}_k,\;
\alpha \sigma_n \mathbf{h}_k^T\mathbf{1}
\right]^T,
\label{eq:init_feature}
\end{equation}
where $\mathbf{h}_k$ is the $k$-th column of $\mathbf{H}$ and $\mathbf{1}$ is the all-ones vector. Stacking all node features yields
\begin{equation}
\mathbf{S}^{(0)}=
[\mathbf{s}_1^{(0)},\mathbf{s}_2^{(0)},\ldots,\mathbf{s}_K^{(0)}]^T
\in\mathbb{R}^{K\times 2}.
\end{equation}

The code supports two graph construction modes:
\begin{equation}
\mathbf{L} =
\begin{cases}
\mathbf{L}_{\mathrm{cep}}, & \texttt{graph\_mode = cep\_l},\\
\mathbf{L}_{\mathrm{corr}}, & \texttt{graph\_mode = channel\_corr},
\end{cases}
\label{eq:graph_mode}
\end{equation}
where $\mathbf{L}_{\mathrm{corr}}$ is a row-normalized channel-correlation graph generated from $\mathbf{H}$ with the diagonal removed. During neural processing, the dense matrix $\mathbf{L}$ is converted to sparse edge indices and edge weights by extracting all nonzero entries, which matches the implementation of \texttt{dense\_batch\_to\_sparse()}.

\section{The Proposed D-EPNet Detector}

\subsection{Overall Iterative Structure}

The proposed detector combines an EP module and a DIFFormer-GRU refinement module. Let $T$ denote the total number of EP iterations. In the code, this number is controlled by \texttt{ep\_iter}, while the damping factor is denoted by \texttt{ep\_beta}. For each iteration $t=1,2,\ldots,T$, the detector performs:
\begin{enumerate}
\item EP-based Gaussian approximation and cavity-statistic extraction;
\item graph neural refinement using DIFFormer and a recurrent hidden state;
\item posterior classification over the alphabet $\mathcal{X}$ and EP site update.
\end{enumerate}

The outer loop is implemented in the \texttt{ChebNet.forward()} function:
\begin{equation}
\big(a_{\mathrm{cav}}^{(t)}, v_{\mathrm{cav}}^{(t)}, \bm{\Lambda}^{(t)}, \bm{\gamma}^{(t)}\big)
=
\mathrm{EP}\!\left(
\mathbf{H}, \mathbf{y}, \sigma_n^2,
\bm{\Lambda}^{(t-1)}, \bm{\gamma}^{(t-1)},
\mathbf{p}^{(t-1)}
\right),
\end{equation}
followed by
\begin{equation}
\big(\mathbf{p}^{(t)}, \mathbf{R}_{\mathrm{gru}}^{(t)}, \mathbf{G}^{(t)}\big)
=
\Phi_{\mathrm{DIF}}\!\left(
\mathbf{S}^{(0)}, \mathbf{L}, \sigma_n^2,
a_{\mathrm{cav}}^{(t)}, v_{\mathrm{cav}}^{(t)},
\mathbf{R}_{\mathrm{gru}}^{(t-1)}, \mathbf{G}^{(t-1)}
\right),
\label{eq:overall_phi}
\end{equation}
where $\mathbf{p}^{(t)}$ denotes the symbol logits generated by the DIFFormer-based refinement network, $\mathbf{R}_{\mathrm{gru}}^{(t)}$ is the feature tensor fed to the next iteration, and $\mathbf{G}^{(t)}$ is the recurrent hidden state.

\subsection{EP Observation Module}

EP approximates the discrete prior in \eqref{eq:posterior} by a Gaussian site function. At iteration $t$, the approximate posterior is
\begin{equation}
q^{(t)}(\mathbf{x}|\mathbf{y})
\propto
p(\mathbf{y}|\mathbf{x})\chi^{(t)}(\mathbf{x}),
\label{eq:q_t}
\end{equation}
with
\begin{equation}
\chi^{(t)}(\mathbf{x})
\propto
\mathcal{N}
\left(
\mathbf{x}:
(\bm{\Lambda}^{(t-1)})^{-1}\bm{\gamma}^{(t-1)},
(\bm{\Lambda}^{(t-1)})^{-1}
\right),
\label{eq:site_function}
\end{equation}
where
\begin{equation}
\bm{\Lambda}^{(t-1)} = \mathrm{diag}(\lambda_1^{(t-1)},\ldots,\lambda_K^{(t-1)})
\end{equation}
and
\begin{equation}
\bm{\gamma}^{(t-1)} = [\gamma_1^{(t-1)},\ldots,\gamma_K^{(t-1)}]^T.
\end{equation}
The implementation initializes
\begin{equation}
\bm{\gamma}^{(0)}=\mathbf{0},
\qquad
\bm{\Lambda}^{(0)}=\frac{3}{4^k-1}\mathbf{I},
\label{eq:init_ep}
\end{equation}
which matches the code line
\texttt{lamda = 3.0 / (4\string**k - 1)}.

The Gaussian approximation can be written as
\begin{equation}
q^{(t)}(\mathbf{x}|\mathbf{y})
=
\mathcal{N}
(\mathbf{x}:\bm{\mu}^{(t)},\bm{\Sigma}^{(t)}),
\end{equation}
where
\begin{equation}
\bm{\Sigma}^{(t)}=
\left(
\sigma_n^{-2}\mathbf{H}^T\mathbf{H}
+
\bm{\Lambda}^{(t-1)}
\right)^{-1},
\label{eq:sigma_t}
\end{equation}
\begin{equation}
\bm{\mu}^{(t)}=
\bm{\Sigma}^{(t)}
\left(
\sigma_n^{-2}\mathbf{H}^T\mathbf{y}
+
\bm{\gamma}^{(t-1)}
\right).
\label{eq:mu_t}
\end{equation}

Let $\mu_k^{(t)}$ and $\Sigma_k^{(t)}$ denote the $k$-th mean and diagonal variance, respectively. The symbol-wise cavity distribution is
\begin{equation}
v_{\mathrm{cav},k}^{(t)}
=
\frac{\Sigma_k^{(t)}}
{1-\Sigma_k^{(t)}\lambda_k^{(t-1)}},
\label{eq:vcav}
\end{equation}
\begin{equation}
a_{\mathrm{cav},k}^{(t)}
=
v_{\mathrm{cav},k}^{(t)}
\left(
\frac{\mu_k^{(t)}}{\Sigma_k^{(t)}}-\gamma_k^{(t-1)}
\right),
\label{eq:acav}
\end{equation}
which gives
\begin{equation}
q_{\backslash k}^{(t)}(x_k)
=
\mathcal{N}
\!\left(
x_k:
a_{\mathrm{cav},k}^{(t)},
v_{\mathrm{cav},k}^{(t)}
\right).
\label{eq:cavity}
\end{equation}

In the implementation, these cavity statistics are returned as \texttt{mean\_ab} and \texttt{var\_ab}, and they are then passed to the neural refinement block instead of being used alone.

\subsection{DIFFormer-GRU Refinement Module}

\subsubsection{Iterative Feature Initialization}

The network \texttt{sp\_GNN\_Difformer} does not rebuild node embeddings from scratch at every iteration. Instead, the initial node feature at iteration $t$ is
\begin{equation}
\mathbf{U}^{(t)}=
\begin{cases}
\mathrm{FC}_{\mathrm{in}}\!\left(\mathbf{S}^{(0)}\right), & t=1,\\
\mathbf{R}_{\mathrm{gru}}^{(t-1)}, & t>1,
\end{cases}
\label{eq:iterative_init}
\end{equation}
where $\mathrm{FC}_{\mathrm{in}}(\cdot)$ is the linear layer \texttt{fc1a}. Thus, the first iteration uses the handcrafted preprocessing feature, while later iterations reuse the recurrent readout from the previous step.

This tensor is then processed by a three-layer multilayer perceptron:
\begin{equation}
\mathbf{F}^{(t)}=
\rho\!\left(
\mathbf{W}_{2c}
\rho\!\left(
\mathcal{D}
\left(
\mathbf{W}_{2b}
\rho\!\left(
\mathcal{D}
\left(
\mathbf{W}_{2a}\mathbf{U}^{(t)}
\right)
\right)
\right)
\right)
\right),
\label{eq:mlp_pre}
\end{equation}
where $\rho(\cdot)$ denotes ReLU and $\mathcal{D}(\cdot)$ denotes dropout. This matches the layers \texttt{fc2a}, \texttt{fc2b}, and \texttt{fc2c} in the code.

\subsubsection{DIFFormer Block}

The processed feature matrix $\mathbf{F}^{(t)}$ is flattened across the batch dimension, while the graph matrix $\mathbf{L}$ is converted into sparse edges $(\mathcal{E}, \mathbf{w})$. The DIFFormer first applies an input projection:
\begin{equation}
\mathbf{X}^{(0)}=
\mathcal{D}
\left(
\rho
\left(
\mathrm{LN}\!\left(
\mathbf{W}_{\mathrm{in}}\mathbf{F}^{(t)}
\right)
\right)
\right),
\label{eq:input_proj}
\end{equation}
where $\mathrm{LN}(\cdot)$ is LayerNorm.

For the $\ell$-th DIFFormer layer and the $h$-th attention head, query, key, and value vectors are
\begin{equation}
\mathbf{q}_{i,h}^{(\ell)} = \mathbf{W}_{Q,h}^{(\ell)}\mathbf{x}_i^{(\ell)},
\quad
\mathbf{k}_{i,h}^{(\ell)} = \mathbf{W}_{K,h}^{(\ell)}\mathbf{x}_i^{(\ell)},
\quad
\mathbf{v}_{i,h}^{(\ell)} = \mathbf{W}_{V,h}^{(\ell)}\mathbf{x}_i^{(\ell)}.
\label{eq:qkv}
\end{equation}
The code uses the \texttt{simple} kernel, which first performs $\ell_2$ normalization:
\begin{equation}
\bar{\mathbf{q}}_{i,h}^{(\ell)}
=
\frac{\mathbf{q}_{i,h}^{(\ell)}}
{\|\mathbf{q}_{i,h}^{(\ell)}\|_2+\epsilon},
\qquad
\bar{\mathbf{k}}_{i,h}^{(\ell)}
=
\frac{\mathbf{k}_{i,h}^{(\ell)}}
{\|\mathbf{k}_{i,h}^{(\ell)}\|_2+\epsilon}.
\label{eq:norm_qk}
\end{equation}
The global diffusion-attention output can be summarized as
\begin{equation}
\mathbf{a}_{i,h}^{(\ell)}
=
\frac{
\bar{\mathbf{q}}_{i,h}^{(\ell)}
\left(
\sum_{j=1}^{K}
\bar{\mathbf{k}}_{j,h}^{(\ell)}
(\mathbf{v}_{j,h}^{(\ell)})^T
\right)
+
\sum_{j=1}^{K}\mathbf{v}_{j,h}^{(\ell)}
}{
(\bar{\mathbf{q}}_{i,h}^{(\ell)})^T
\left(
\sum_{j=1}^{K}\bar{\mathbf{k}}_{j,h}^{(\ell)}
\right)
+ K
},
\label{eq:linear_attention}
\end{equation}
which reflects the linear-time attention implemented in \texttt{full\_attention()}.

In parallel, the graph branch performs normalized weighted aggregation:
\begin{equation}
\mathbf{g}_{i,h}^{(\ell)}
=
\sum_{j\in\mathcal{N}(i)}
\frac{w_{ij}}{\sqrt{d_i d_j}}
\mathbf{v}_{j,h}^{(\ell)},
\label{eq:gcn_branch}
\end{equation}
where $w_{ij}$ is the edge weight extracted from $\mathbf{L}$. With the default setting \texttt{graph\_weight = -1}, the attention branch and graph branch are added directly. After averaging over all heads, the intermediate output is
\begin{equation}
\mathbf{t}_i^{(\ell)}
=
\frac{1}{H_{\mathrm{att}}}
\sum_{h=1}^{H_{\mathrm{att}}}
\left(
\mathbf{a}_{i,h}^{(\ell)} + \mathbf{g}_{i,h}^{(\ell)}
\right).
\label{eq:transconv_out}
\end{equation}

The residual update used in \texttt{DIFFormer\_v2} is
\begin{equation}
\mathbf{x}_i^{(\ell+1)}
=
\rho
\left(
\mathcal{D}
\left(
\mathrm{LN}
\left(
\alpha_{\mathrm{res}}\mathbf{t}_i^{(\ell)}
+
(1-\alpha_{\mathrm{res}})\mathbf{x}_i^{(\ell)}
\right)
\right)
\right),
\label{eq:residual_update}
\end{equation}
where $\alpha_{\mathrm{res}}=0.5$ by default in the provided implementation. After $L$ DIFFormer layers, an output linear layer maps the hidden representation back to the state dimension $s_u$.

\subsubsection{Fusion with EP Statistics and GRU Memory}

The DIFFormer output alone is not used as the final prediction. Instead, it is fused with the EP cavity statistics:
\begin{equation}
\mathbf{c}_k^{(t)}
=
\left[
\mathbf{m}_k^{(t)};
a_{\mathrm{cav},k}^{(t)};
v_{\mathrm{cav},k}^{(t)}
\right],
\label{eq:gru_input}
\end{equation}
where $\mathbf{m}_k^{(t)}$ is the DIFFormer message for the $k$-th symbol. This concatenated vector is fed into a time-distributed GRU cell:
\begin{equation}
\mathbf{g}_k^{(t)}
=
\mathrm{GRUCell}
\left(
\mathbf{c}_k^{(t)},
\mathbf{g}_k^{(t-1)}
\right).
\label{eq:gru_update}
\end{equation}
The hidden state is then mapped back to the state space by
\begin{equation}
\mathbf{r}_k^{(t)}=\mathbf{W}_{4}\mathbf{g}_k^{(t)}.
\label{eq:read_gru}
\end{equation}
The tensor $\mathbf{r}_k^{(t)}$ is reused in \eqref{eq:iterative_init} at the next iteration, which means the neural module carries memory across EP rounds.

Finally, a three-layer readout MLP produces the symbol logits:
\begin{equation}
\mathbf{o}_k^{(t)}
=
\mathbf{W}_{3c}
\mathbf{W}_{3b}
\mathcal{D}\!\left(
\mathbf{W}_{3a}
\mathbf{r}_k^{(t)}
\right)
\end{equation}
Equivalently, one may write
\begin{equation}
\mathbf{u}_k^{(t)}=\mathcal{D}\!\left(\mathbf{W}_{3a}\mathbf{r}_k^{(t)}\right),\qquad
\mathbf{v}_k^{(t)}=\mathbf{W}_{3b}\mathbf{u}_k^{(t)},\qquad
\mathbf{o}_k^{(t)}=\mathbf{W}_{3c}\mathbf{v}_k^{(t)}.
\label{eq:readout}
\end{equation}
In implementation, dropout is applied only after the first readout layer. The class posterior is
\begin{equation}
p_{\mathrm{DIF}}^{(t)}(x_k=a_m|\mathbf{y},\mathbf{H})
=
\frac{\exp(o_{k,m}^{(t)})}
{\sum_{j=1}^{|\mathcal{X}|}\exp(o_{k,j}^{(t)})},
\qquad a_m\in\mathcal{X}.
\label{eq:softmax}
\end{equation}

\subsection{EP Site Update and Final Detection}

Using the class posterior in \eqref{eq:softmax}, the refined mean and variance are
\begin{equation}
\hat{x}_k^{(t)}
=
\sum_{a\in\mathcal{X}}
a\,p_{\mathrm{DIF}}^{(t)}(x_k=a|\mathbf{y},\mathbf{H}),
\label{eq:xhat}
\end{equation}
\begin{equation}
\hat{v}_k^{(t)}
=
\sum_{a\in\mathcal{X}}
\left(a-\hat{x}_k^{(t)}\right)^2
p_{\mathrm{DIF}}^{(t)}(x_k=a|\mathbf{y},\mathbf{H}).
\label{eq:vhat}
\end{equation}
Let
\begin{equation}
\hat{\mathbf{x}}^{(t)}
=
[\hat{x}_1^{(t)},\ldots,\hat{x}_K^{(t)}]^T,
\qquad
\hat{\mathbf{V}}^{(t)}
=
\mathrm{diag}(\hat{v}_1^{(t)},\ldots,\hat{v}_K^{(t)}).
\end{equation}
The EP site parameters are updated as
\begin{equation}
\bm{\Lambda}^{(t)}
=
(\hat{\mathbf{V}}^{(t)})^{-1}
-
(\mathbf{V}_{\mathrm{cav}}^{(t)})^{-1},
\label{eq:lambda_update}
\end{equation}
\begin{equation}
\bm{\gamma}^{(t)}
=
(\hat{\mathbf{V}}^{(t)})^{-1}\hat{\mathbf{x}}^{(t)}
-
(\mathbf{V}_{\mathrm{cav}}^{(t)})^{-1}\mathbf{a}_{\mathrm{cav}}^{(t)},
\label{eq:gamma_update}
\end{equation}
with
\begin{equation}
\mathbf{a}_{\mathrm{cav}}^{(t)}
=
[a_{\mathrm{cav},1}^{(t)},\ldots,a_{\mathrm{cav},K}^{(t)}]^T
\end{equation}
and
\begin{equation}
\mathbf{V}_{\mathrm{cav}}^{(t)}
=
\mathrm{diag}(v_{\mathrm{cav},1}^{(t)},\ldots,v_{\mathrm{cav},K}^{(t)}).
\end{equation}

To improve numerical stability, negative precision terms are clipped by reusing the previous EP values. Moreover, the code applies damping inside \texttt{performEP()} with coefficient $\beta$, which corresponds to the command-line argument \texttt{ep\_beta}. At the final iteration, the hard decision is obtained by nearest-neighbor slicing:
\begin{equation}
\hat{x}_k
=
\arg\min_{a\in\mathcal{X}}
|\hat{x}_k^{(T)}-a|.
\label{eq:hard_decision}
\end{equation}

\section{Training Objective and Implementation Details}

The neural module outputs a logit vector for each real-valued symbol. Let $c_{b,k}$ denote the ground-truth class index of the $k$-th symbol in the $b$-th sample. The code trains the network by cross-entropy:
\begin{equation}
\mathcal{L}
=
\sum_{b=1}^{B}\sum_{k=1}^{K}
\mathrm{CE}\!\left(
\mathbf{o}_{b,k}, c_{b,k}
\right),
\label{eq:loss}
\end{equation}
where $\mathbf{o}_{b,k}$ is the output logit vector before softmax.

For performance evaluation, the soft symbol estimate is computed from the posterior probabilities as a constellation-weighted sum, and then mapped back to the nearest discrete symbol index. The training script uses Adam for optimization and \texttt{ReduceLROnPlateau} for learning-rate scheduling. The number of DIFFormer layers, dropout ratio, EP iteration count, damping factor, and graph construction mode are all exposed as command-line arguments, which makes the detector architecture configurable without changing the inference pipeline.

\section{Conclusion}

This draft reformulates the method description according to the current implementation. In particular, the proposed detector should be described as an EP-guided DIFFormer-GRU architecture rather than a Chebyshev graph detector. The key implementation characteristics are: 1) EP-based cavity inference, 2) optional CEP or channel-correlation graph construction, 3) linear-time diffusion attention combined with graph aggregation, and 4) recurrent feature reuse across EP iterations.

\begin{thebibliography}{00}
\bibitem{b1} T. P. Minka, ``Expectation propagation for approximate Bayesian inference,'' in \emph{Proc. Conf. Uncertainty in Artificial Intelligence}, 2001, pp. 362--369.
\bibitem{b2} A. Vaswani \emph{et al.}, ``Attention is all you need,'' in \emph{Advances in Neural Information Processing Systems}, 2017, pp. 5998--6008.
\bibitem{b3} IEEE, ``IEEE conference templates,'' Accessed: Apr. 14, 2026. [Online]. Available: \url{https://www.ieee.org/conferences/publishing/templates.html}
\end{thebibliography}

\end{document}
